\subsection{由某些子群的 Haar 测度构造群的 Haar 测度}

设 \(\mathbf{G}\) 是局部紧群，\(\mathbf{X}\) 和 \(\mathbf{Y}\) 是 \(\mathbf{G}\) 的两个闭子群，且 \(\Omega = \mathbf{XY}\) 包含一个邻域 \(\mathbf{U}\)，其中包含 \(e\)。则 \(\Omega\) 是 \(\mathbf{G}\) 中的开集；事实上，对于任意 \(x_0 \in \mathbf{X}\) 和 \(y_0 \in \mathbf{Y}\)，有 \(\mathbf{XY} = (x_0\mathbf{X})(\mathbf{Y}y_0) \supset x_0\mathbf{U}y_0\)，而 \(x_{0}Uy_{0}\) 是 \(x_{0}y_{0}\) 的邻域；因此 \(\Omega\) 是其每一点的邻域。

*当 G 是李群且其李代数为 g 时，如果 X、Y 所对应的子代数之和为 g，则关于 X、Y 的条件得到满足。*

群 \(X \times Y\) 按作用律 \((x, y) \cdot s = xsy^{-1}\)（\(x \in X, y \in Y, s \in G\)）在 G 上连续地从左作用。令 \(Z = X \cap Y\)。\(X \times Y\) 中 e 的稳定子是子群 \(Z_{0}\)（它属于 \(X \times Y\)），它由 \((z, z)\)（其中 \(z \in Z\)）这些对组成，并且是与 Z 典范同构的子群。因此，集合 \(\Omega\) 可识别为左陪集齐性空间 \((X \times Y)/Z_{0}\)；更准确地说，映射 \((x, y) \mapsto xy^{-1}\) 从 \(X \times Y\) 到 \(\Omega\)，通过取商定义了从 \((X \times Y)/Z_{0}\) 到 \(\Omega\) 的连续双射。我们假设该映射是同胚。（当 G 在无穷远处可数时尤其如此：参见 App. I。）

命题 13. --- 进一步假设 Z 是紧的。设 \(\mu_{\mathrm{G}}\)、\(\mu_{\mathrm{X}}\)、\(\mu_{\mathrm{Y}}\) 分别是 G、X、Y 上的左 Haar 测度，\(\Lambda\) 是 \(\Delta_{\mathrm{G}}\) 在 Y 上的限制。则限制 \(\mu\)（即 \(\mu_{\mathrm{G}}\) 在 \(\Omega\) 上的限制），除一个常数因子外，是 \(\mu_{\mathrm{X}} \otimes (\Lambda^{-1} \cdot \mu_{\mathrm{Y}})\) 在映射 \((x, y) \mapsto xy^{-1}\) 从 X × Y 到 \(\Omega\) 下的像（该映射是 逆紧的）。

对于 \(x\in \mathbf{X},y\in \mathbf{Y}\)，

\[
\boldsymbol {\gamma} \big ((x, y) \big) \mu = \boldsymbol {\delta} (y) \boldsymbol {\gamma} (x) \mu = \Delta_ {G} (y) \mu .
\]

将 \(\Omega\) 与齐性空间 \((\mathrm{X} \times \mathrm{Y}) / \mathrm{Z}_0\) 识别，并在 \(\mathbf{Z}_0\) 上选取适当的 Haar 测度，可见 \(\mu^\sharp\) 是 \(\mathrm{X} \times \mathrm{Y}\) 的左 Haar 测度（即 \(\mu_{\mathrm{X}} \otimes \mu_{\mathrm{Y}}\)）乘以函数 \((x, y) \mapsto \Delta_{\mathrm{G}}(y)^{-1}\) 所得的测度（No.~6, 引理 6）。另一方面，\(\mu\) 除一个常数因子外，是 \(\mu^\sharp\) 在从 \(\mathrm{X} \times \mathrm{Y}\) 到 \(\Omega\) 的典范映射下的像（No.~3, 注记 2）。

推论。--- 设 f 是定义在 \(\Omega\) 上、取值于 Banach 空间或 \(\overline{R}\) 的函数。f 关于 \(\mu\) 可积的充要条件是，函数 \((x,y)\mapsto\mathbf{f}(xy)\Delta_{\mathrm{G}}(y)\Delta_{\mathrm{Y}}(y)^{-1}\) 关于 \((\mu_{\mathrm{X}}\otimes\mu_{\mathrm{Y}})\) 可积；在此情形下

\[
\int_ {\Omega} \mathbf {f} (\omega) d \mu (\omega) = a \iint_ {\mathrm{X} \times \mathrm{Y}} \mathbf {f} (x y) \Delta_ {\mathrm{G}} (y) \Delta_ {\mathrm{Y}} (y) ^ {- 1} d \mu_ {\mathrm{X}} (x) d \mu_ {\mathrm{Y}} (y),\tag{23}
\]

其中 \(a\) 是一个 \(>0\) 且与 \(\mathbf{f}\) 无关的常数。

由 Prop. 13 和 Ch. V, §4, No.~4, Th. 2，f 关于 \(\mu\) 可积的充要条件是，函数 \((x,y)\mapsto\mathbf{f}(xy^{-1})\) 关于 \(\mu_{\mathrm{X}}\otimes(\Lambda^{-1}\cdot\mu_{\mathrm{Y}})\) 可积；或者等价地，函数 \((x,y)\mapsto\mathbf{f}(xy^{-1})\Delta_{\mathrm{G}}(y)^{-1}\) 关于 \(\mu_{X}\otimes\mu_{Y}\) 可积；或者等价地，函数 \((x,y)\mapsto\mathbf{f}(xy)\Delta_{\mathrm{G}}(y)\Delta_{\mathrm{Y}}(y)^{-1}\) 关于 \(\mu_{X}\otimes\mu_{Y}\) 可积。公式 (23) 由类似论证得出。

命题 14. --- 假设 Prop. 13 的条件成立，并且进一步假设 Y 是正规的。

\begin{enumerate}
\def\labelenumi{\alph{enumi})}
\item
  \(\mu_{\mathrm{G}}\) 在 \(\Omega\) 上的限制，除一个常数因子外，是 \(\mu_{\mathrm{X}} \otimes \mu_{\mathrm{Y}}\) 在映射 \((x,y) \mapsto xy\) 从 \(\mathrm{X} \times \mathrm{Y}\) 到 \(\Omega\) 下的像。
\item
  对于 \(x \in \mathbf{X}\) 和 \(y \in \mathbf{Y}\)，
\end{enumerate}

\[
\Delta_ {\mathrm{G}} (x y) = \Delta_ {\mathrm{X}} (x) \Delta_ {\mathrm{Y}} (y) \bmod (i _ {x}),
\]

其中 \(i_x\) 表示 Y 的自同构 \(v \mapsto x^{-1}vx\)。

由 Prop. 10 b)，\(\Delta_{\mathrm{G}} = \Delta_{\mathrm{Y}}\) 在 Y 上成立，因此 a) 由 (23) 得出。取 \(x_0 \in \mathbf{X}\)、\(y_0 \in \mathbf{Y}\)。记 \(p\) 为映射 \((x, y) \mapsto xy\)，从 \(\mathbf{X} \times \mathbf{Y}\) 到 \(\Omega\)。由于

\[
x y \left(x _ {0} y _ {0}\right) ^ {- 1} = x x _ {0} ^ {- 1} \left(x _ {0} y y _ {0} ^ {- 1} x _ {0} ^ {- 1}\right) = x x _ {0} ^ {- 1} i _ {x _ {0} ^ {- 1}} \left(y y _ {0} ^ {- 1}\right),
\]

我们有

\[
\begin{array}{r l} \Delta_ {\mathrm{G}} (x _ {0} y _ {0}) p (\mu_ {\mathrm{X}} \otimes \mu_ {\mathrm{Y}}) & = \pmb {\delta} (x _ {0} y _ {0}) p (\mu_ {\mathrm{X}} \otimes \mu_ {\mathrm{Y}}) \\ & = p \big (\pmb {\delta} (x _ {0}) \mu_ {\mathrm{X}} \otimes i _ {x _ {0} ^ {- 1}} \pmb {\delta} (y _ {0}) \mu_ {\mathrm{Y}} \big) \\ & = p \big (\Delta_ {\mathrm{X}} (x _ {0}) \mu_ {\mathrm{X}} \otimes \Delta_ {\mathrm{Y}} (y _ {0}) (\bmod i _ {x _ {0}}) \mu_ {\mathrm{Y}} \big) \\ & = \Delta_ {\mathrm{X}} (x _ {0}) \Delta_ {\mathrm{Y}} (y _ {0}) (\bmod i _ {x _ {0}}) p (\mu_ {\mathrm{X}} \otimes \mu_ {\mathrm{Y}}), \end{array}
\]

从而得到 b)。

注记。--- 当 G 是 X 与 Y 的拓扑半直积时，Prop. 14 特别适用（GT, III, §2, No.~10）。此时 \(Z = \{e\}\) 且 \(\Omega = G\)。由于有 \(yx = x i_{x}(y)\)，对 \(x \in X\)、\(y \in Y\) 成立，测度 \(\mu_{G}\) 也等于 \((\text{mod } i_{x})\mu_{X} \otimes \mu_{Y}\) 在映射 \((x, y) \mapsto yx\) 从 \(X \times Y\) 到 G 下的像，差一个常数因子。

